\section{Integrated Discussion}
\label{sec:discussion}

\subsection{Model Hierarchy and Question Matching}
\label{subsec:discussion-model-hierarchy-question-matching}

The first research question asked how blood-flow model classes are organized
and what distinguishes them. The review answers that question by treating the
hierarchy as a set of state choices rather than a ranking of accuracy. A 3D CFD
model retains local velocity, pressure, stress, and geometry. FSI adds wall
motion and interface work. Axisymmetric and 2D models keep selected resolved
structure under stronger geometric assumptions. Distributed 1D models keep
axial area, flow, mean velocity, and pressure or wall-law variables. 0D
networks keep pressure--flow relations without local stenosis geometry.
Multiscale models couple these tiers, and learned models approximate maps whose
meaning depends on the tier and data used to train them.

The practical consequence is that a simulation should be read from the question
backward. If the question concerns wall traction, a reduced area--flow state is
not the native object. If the question concerns network pulse propagation, a
fully resolved local geometry may be unnecessary without appropriate boundary
data. If the question concerns agreement between tiers, an observation operator
is part of the mathematical claim. The idealized stenosis example illustrates
this rule by placing a distributed 1D state beside a resolved velocity field
only after both are mapped to common flow and section-mean velocity
observables.

\subsection{Observables and Closure Choices}
\label{subsec:discussion-closures-interpretability}

The second research question asked how closure choices, data, numerical
discretization, and output operators affect interpretation. The review shows
that rheology, wall mechanics, geometry, boundary data, and observables are not
secondary labels. They define the model. A Newtonian stress law and a
generalized-Newtonian stress law can change viscous and shear interpretation. A
fixed wall, reduced pressure--area law, and resolved FSI wall represent
different mechanical states. A Windkessel outlet, pressure outlet, traction
condition, and characteristic boundary rule impose different downstream
information. A section-mean velocity, wall shear stress, pressure ratio, and
flow rate are different observables.

The case study makes this interpretive contract explicit without turning it
into a clinical validation exercise. It uses a smooth stenosis geometry,
reduced area--flow variables, a physical-to-solver coordinate map, a
pressure--area wall law, parabolic-profile closure, Newtonian viscosity for the
principal comparison, and a finite-volume realization. That list is not a
software inventory. It is the mathematical information needed before a
computed quantity can be read.

\subsection{Numerical Evidence and Cross-Dimensional Interpretation}
\label{subsec:discussion-numerical-evidence-observation}

The numerical review separates implementation verification, equilibrium
preservation, sensitivity, cross-model comparison, reproducibility, and
validation. Manufactured solutions check whether a declared operator reproduces
a known smooth state. Grid and timestep studies check resolution sensitivity.
Well-balanced tests ask whether a discrete method preserves a meaningful
equilibrium. Benchmarks and open workflows support reproducibility. External
validation requires an accepted target for the specific observable.

The worked stenosis example demonstrates why these evidence categories should
not be collapsed. The manufactured-solution record supports the declared
forced operator. The geometry-rest test asks a different question about
flux-source cancellation in a spatially varying vessel. The
plane--tetrahedron quadrature operator asks a third question: how a resolved 3D
velocity field is reduced to section area, physical flow, and area-mean
velocity. Agreement or limitation in one category does not replace evidence in
another. The recorded 1D--3D comparison is therefore a descriptive
cross-model discrepancy under a declared observation operator, not a validation
error.

\subsection{What the Case Study Contributes}
\label{subsec:discussion-case-study-contribution}

The idealized stenosis computation contributes one controlled application of the
review framework: a methodology for interpreting a reduced-model comparison,
not a predictive stenosis result. It shows how a controlled geometry can isolate
model-reduction and numerical-method questions; how solver coordinates must be
mapped back to physical area and flow; how boundary, wall, profile, and
rheology closures enter the model contract; how equilibrium preservation
becomes a discrete property; and how resolved and reduced models require a
shared observation map before their outputs can be compared.

The same example also shows what remains open. Stronger production comparison
claims would require a repaired or explicitly well-balanced discretization if
the rest-state test remains limiting, matched resolved-data metadata for wall,
boundary, material, geometry state, and time, and production sensitivity
studies for the reported comparison outputs. Those are not cosmetic additions.
They are the evidence needed to move from a bounded worked example to a broader
predictive or validation claim.

\section{Conclusion}
\label{sec:conclusions-limitations}

This report provides a review-first structure for interpreting mathematical
and numerical blood-flow models. It connects continuum balance laws, CFD and
FSI models, reduced and lumped hemodynamic models, multiscale coupling,
constitutive and wall closures, boundary conditions, numerical methods,
verification, validation, and observation operators.

The central conclusion is that a hemodynamic result is meaningful only relative
to the equations, retained state, closures, boundary data, discretization,
evidence standard, and observable that produced it. The literature synthesis
shows that fidelity, cost, identifiability, reproducibility, and validation
status must be matched to the question being asked. The idealized stenosis case
study supplies one bounded example of that standard: it specifies the reduced
model, states the relevant closure choices and numerical evidence record, tests an
equilibrium property, and compares reduced and resolved velocity information
through a declared quadrature operator.

The report does not claim clinical validation of the case study or universal
superiority of a model tier. Its contribution is methodological: it organizes
the blood-flow modeling landscape around model hierarchy, closure choices,
observables, numerical evidence, and interpretation limits. Future work should
strengthen the same chain where needed: well-balanced discretization for the
geometry-rest state, matched boundary and wall metadata for resolved
comparisons, and declared validation observables before predictive or clinical
claims are made.
